\newcommand{\bih}{{\boldmath${h}$}}

\newcommand{\biP}{{\boldmath${P}$}}

\newcommand{\bip}{{\boldmath${p}$}}



\def\stat{zabez}



\def\tit{К ЗАДАЧЕ АНАЛИЗА ВЛОЖИМОСТИ ПОДСЛОВ В~ЗАГОЛОВКИ ПАКЕТОВ ДАННЫХ}



\def\titkol{К задаче анализа вложимости подслов в~заголовки пакетов данных} 



\def\autkol{М.\,И.~Забежайло}



\def\aut{М.\,И.~Забежайло$^1$}



\titel{\tit}{\aut}{\autkol}{\titkol}



%{\renewcommand{\thefootnote}{\fnsymbol{footnote}}\footnotetext[1]

%{Работа выполнена при поддержке РФФИ (грант 11-07-00112).}}



\renewcommand{\thefootnote}{\arabic{footnote}}

\footnotetext[1]{Центр прикладных исследований компьютерных сетей, Сколково, MZabezhailo@arccn.ru}





\Abst{Обсуждаются возможности использования технологий про\-грам\-мно-кон\-фи\-гу\-ри\-ру\-емых сетей 

(ПКС-тех\-но\-ло\-гий) в оптимизации управ\-ле\-ния сетевым трафиком. Предложена алгебраическая 

формализация задачи вложимости подслов в слова заданного алфавита. Сформулированы условия 

и предложены алгоритмы, позволяющие оптимизировать процедуры проверки вложимости 

подслов ограниченной длины в строки таблиц коммутации больших размеров. Показаны 

возможности дополнительного ускорения обработки пакетов данных в компьютерных сетях с 

использованием ПКС-тех\-но\-логий.}



\KW{программно-конфигурируемые сети; анализ заголовков пакетов данных в компьютерных 

сетях; математические методы анализа данных}



\vskip 14pt plus 9pt minus 6pt





      \thispagestyle{headings}



      

            \label{st\stat}

            

            \vspace*{-12pt}



\section{Введение}



     Одними из ключевых элементов современного комплекса технологий передачи 

данных в компьютерных сетях являются процедуры проверки вложимости символьных 

наборов заданного вида в те или иные объекты. Среди них такие процедуры, как проверка 

встречаемости заголовка конкретного пакета данных в соответствующей таблице 

коммутации, проверка встречаемости подслов заданного вида в передаваемых данных 

и~т.\,п. (см., например,~[1]). 

     

     В частности, вложимость подслов заданного вида в заголовок передаваемого по сети 

пакета данных может определять конкретный вид набора процедур обработки этого 

пакета в соответствующих узлах сети. При этом быстрота выделения соответствующего 

набора процедур обработки~--- существенный элемент, вли\-яющий на скорость 

прохождения данных в сети. Другой не менее чувствительной (в части обеспечения 

нормального режима функционирования компьютерных систем и сетей) процедурой 

является быстрое опознание (выделение в теле передаваемого пакета данных) так 

называемых сигнатур вредоносного программного обеспечения (ПО). Сегодня при количестве 

известных компьютерных вирусов, оцениваемом в десятки тысяч единиц, объем базы 

характеризующих их сигнатур оценивается уже в сотни тысяч единиц, а характерные 

времена адекватной (обеспечивающих защиту) реакции соответствующих антивирусных 

служб на выявление того или иного вредоносного ПО~--- десятки часов (в то время как 

оценки характерного времени возможного заражения существенных по объемам 

фрагментов сети Интернет оказываются на порядок меньше~--- десятки минут). Таким 

образом, в ситуации, когда процедуры обсуждаемого типа используются постоянно и 

многократно, практически любые наработки по оптимизации реализуемых в них 

алгоритмов могут рассматриваться как актуальные. 

     

     Массовый характер использования таких процедур выдвигает достаточно жесткие 

требования к их организации: проблема эффективности (оцениваемой прежде всего с 

точки зрения сложности вычислений, см., например,~[2]) переходит в центр внимания 

разработчиков и эксплуатационных служб.

     

     Новые возможности в организации обсуждаемых процедур открывает использование 

так называемых Open Flow технологий (см., например,~[3--7]). Для обозначения этих 

технологий также используется термин Software Defined Networks (SDN)~--- 

     про\-грам\-мно-кон\-фи\-гу\-ри\-ру\-емые сети. Развиваемый здесь подход 

позволяет разделить уровни собственно трафика и управления трафиком в компьютерной 

сети, предоставляя дополнительные возможности в части последнего вынесением 

значительного объема соответствующих функций на внешний контроллер (отдельный 

вычислительный комплекс, взаимодействующий со стандартными сетевыми 

устройствами~--- коммутаторами, маршрутизаторами и~т.\,п.~--- через специальный 

открытый протокол). Таким образом удается существенно повысить производительность 

(и пропускную способность) коммуникационного оборудования, задействованного на 

уровне сетевого трафика, выводя <<интеллектуальные>> функции анализа данных и 

управления на внешний проб\-лем\-но-ори\-ен\-ти\-ро\-ван\-ный ПКС-кон\-трол\-лер.

     

     Некоторые алгоритмические аспекты распределения в рамках ПКС-под\-хо\-да 

функциональности представленных выше процедур (проверки вложимости символьных 

наборов заданного вида в те или иные объекты) между уровнями исполнения трафика и 

управления трафиком и станут предметом дальнейшего рассмотрения.

     



\section{Основные понятия и обозначения}



     Заголовком пакета (далее~--- заголовком) будем называть булевский вектор 

\bih\;$\hm = \langle \alpha_1, \alpha_2, \ldots , \alpha_n\rangle$ фиксированной длины $n$ 

(где для каждого $i$ из множества $\{1, 2, \ldots , n\}$ значение $\alpha_i\hm\in\{0,\,1\}$). 

Пару векторов $\{\mbox{\bih}_{0i}, \mbox{\bih}_{1i}\}$ таких, что на всех позициях, кроме позиции 

номер~$i$, у них расположены одинаковые значения $\alpha_j \in\{0,\,1\}$, а кроме того, 

на позиции номер~$i$ первый из них имеет~0, а второй~--- 1, будем обозначать 

посредством~${h}_{xi}$. Наконец, $\mbox{\bih}_{x\mathbf{I}}$ есть компактная запись для множества всех 

пар векторов вида~$\mbox{\bih}_{xi}$, где $\mathbf{I}$ есть некоторое заданное множество 

значений индексов $i$ в заголовке~$\mbox{\bih}$. Другими словами, в (мета)заголовке $\mbox{\bih}_{x\mathbf{I}}$ 

на всех\footnote{При этом случай, когда на каждой из $n$~позиций вектора~$\mathbf{h}$ 

расположено значение~${x}$, не рассматривается по содержательным 

соображениям.} позициях с номерами~$i$ из множества~$\mathbf{I}$ (и только в них!) 

расположены значения~${x}$: 

     $$

     \mbox{\bih}_{x\mathbf{I}}=\langle \beta_1,\beta_2, \ldots ,\beta_n\rangle\,,

     $$

     где $\beta_j \in\{0,\,1,\,x\}$ и $j\in\{1, 2, \ldots , n\}$.

     

    Рассмотрим алфавит

    $$

    \mathbf{U}^* =\left\{ a_{01},a_{11},a_{02},a_{12}, \ldots , a_{0n},a_{1n}\right\}\,,

    $$

образованный ровно $2n$ различными символами (буквами, образующими). Каж\-до\-му 

заголовку $\mathbf{h}\hm = \langle \alpha_1,\alpha_2,\ldots ,\alpha_n\rangle$ длины~$n$, 

построенному лишь с помощью символов~0 и~1, сопоставим множество $\mathbf{h}\hm = 

\{a_1, a_2, \ldots , a_k\}$ такое, что в качестве элемента (образующей) $a_i$ выбирается 

$a_{0i}$, если $\alpha_i \hm= 0$, и~$a_{1i}$, если $\alpha_i \hm= 1$. Далее, для $\mbox{\bih}_{x\mathbf{I}}\hm = 

\langle\beta_1, \beta_2, \ldots , \beta_n\rangle$ по каждой позиции из 

множества~$\mathbf{I}$ (т.\,е.\ позиции, на которой в $\mbox{\bih}_{x\mathbf{I}}$ находится третье 

значение~--- ${x}$) выберем в соответствующее множество~$\mathbf{h}$ 

одновременно обе образующие~$a_{0i}$ и~$a_{1i}$. Таким образом, появляется 

возможность вза\-им\-но-од\-но\-знач\-но кодировать представленные \textit{кортежами} 

заголовки вида~$\mbox{\bih}$ с помощью множеств вида~$\mathbf{h}$.

     

Будем называть \textit{таблицей коммутации} множество\footnote[2]{Говоря более аккуратно, 

следовало бы 

подмножествам множества~$\mathbf{K}$ заголовков на входе рассматриваемой коммутационной таблицы 

сопоставить номера соответствующих выходных портов (коммутатора, использующего эту 

коммутационную таблицу). Тем не менее, можно (без потери общности), например, считать, что все 

заголовки множества~$\mathbf{K}$ адресуются на один порт, а сводная таблица коммутации формируется 

объединением соответствующих подтаблиц по всем задействованным портам коммутатора.} 

$\mathbf{K}\hm = \{h_1, 

h_2, \ldots , h_m\}$, перечисляющее последовательно заданные заголовки~$h_1, h_2, \ldots , h_m$ 

(переформулированные по представленному выше алгоритму в алфавите~$\mathbf{U}^*$).

    



\section{Некоторые полезные свойства таблиц коммутации}



     Для удобства последующих операций со строками коммутационной таблицы 

определим следующее

     

     \medskip

     

     \noindent

\textbf{Правило.}\ Пусть задан заголовок $\mathbf{h}\hm =\{a_1, a_2, \ldots , a_n\}$, в котором на 

позициях~$i_1, i_2, \ldots , i_s$ ($0 \hm\leq  s \hm<n$) размещаются символы~${x}$ 

(т.\,е.\ в каждой из этих позиций можно разместить как~0, так и~1). Множество $\mathbf{h}^*\hm = 

\{{h}, {h}^{10}, {h}^{11}, \ldots , {h}^{s0}, {h}^{s1}\}$ 

строится так, что для каждого значения~$i$ из имеющегося множества $\{i_1, i_2, \ldots , 

i_s\}$ номеров позиций символа~${x}$ в заголовке~$\mathbf{h}$ каждая пара 

${h}^{i0}, {h}^{i1}$ порождается из исходного заголовка~$\mathbf{h}$ 

заменой~${x}$ на соответствующей позиции один раз на~0, а другой~--- на~1 (см.\ 

таблицу).



%\begin{table}

{\small

\begin{center}

\tabcolsep=4pt

\begin{tabular}{|c|c|c|c|c|c|c|c|c|c|}

\multicolumn{10}{c}{Представления заголовков}\\[2ex]

\hline

& \cellcolor[gray]{.85}$\alpha_1$ & \cellcolor[gray]{.85}$\alpha_2$ & \cellcolor[gray]{.85}$\cdots$ &

\cellcolor[gray]{.85}$\alpha_i$ & \cellcolor[gray]{.85}$\cdots$ & \cellcolor[gray]{.85}$\alpha_j$ &

\cellcolor[gray]{.85}$\cdots$ & 

\cellcolor[gray]{.85}$\alpha_{n-1}$ & \cellcolor[gray]{.85}$\alpha_n$\\

\cline{2-10}

\cline{2-10}

\cline{2-10}

& \cellcolor[gray]{.85}$a_{01}, a_{11}$ &\cellcolor[gray]{.85}$a_{01}, a_{12}$ & 

\cellcolor[gray]{.85}$\cdots$ & \cellcolor[gray]{.85}$a_{0i},a_{1i}$ & 

\cellcolor[gray]{.85}$\cdots $& \cellcolor[gray]{.85}$a_{0j},a_{1j}$ & 

\cellcolor[gray]{.85}$\cdots$ & \cellcolor[gray]{.85}$a_{0\,n-1},a_{1\,n-1}$ & 

\cellcolor[gray]{.85}$a_{0n},a_{1n}$\\

\hline

\cellcolor[gray]{.85}$\mathbf{h}$ & $a_1$ &$a_1$  &$\cdots$ & $x$ & $\cdots$ & $x$ & $\cdots$ & $a_{n-1}$ & 

$a_n$\\

%\hline

\cellcolor[gray]{.85}$\mathbf{h}^{i0}$ & $a_1$ & $a_2$ & $\cdots$ & 0 & $\cdots$ & $x$ & $\cdots$ & 

$a_{n-1}$ & $a_n$\\

%\hline

\cellcolor[gray]{.85}$\mathbf{h}^{i1}$ & $a_1$ & $a_2$ & $\cdots$ & 1 & $\cdots$ & $x$ & $\cdots$ & 

$a_{n-1}$ & $a_n$\\

\cellcolor[gray]{.85}$\cdots$ &$\cdots$ &$\cdots$ &$\cdots$ &$\cdots$ &$\cdots$ &$\cdots$ &$\cdots$ &

$\cdots$ &$\cdots$\\

\cellcolor[gray]{.85}$\mathbf{h}^{j0}$ & $a_1$ & $a_2$ & $\cdots$ & $x$ & $\cdots$ & 0 & $\cdots$ & 

$a_{n-1}$ & $a_n$\\

\cellcolor[gray]{.85}$\mathbf{h}^{j1}$ & $a_1$ & $a_2$ & $\cdots$ & $x$ &$\cdots$ & 1 &  $\cdots$ & 

$a_{n-1}$ & $a_n$\\

\hline

\end{tabular}

\end{center}}

%\end{table}

    

     Результатом применения введенного {правила} к заданной таблице коммутации~$\mathbf{K}$ 

     будет расширенная соответствующим образом таблица $\mathbf{K}^*\hm = \{h_1, h_2, \ldots , 

h_M\}$. 

     

     \medskip

     

     \noindent

     \textbf{Утверждение 1.} \textit{При заданном $n$~--- размерности анализируемых 

заголовков~$\mathbf{h}$~--- построение с применением введенного \textit{правила} по 

исходной таблице $\mathbf{K}\hm = \{h_1, h_2, \ldots , h_m\}$ расширенной таблицы коммутации 

$\mathbf{K}^*\hm = \{h_1, h_2, \ldots , h_M\}$ требует не более чем $O(n)$ операций (т.\,е.\ таблица 

$\mathbf{K}^*$ строится по таблице~$\mathbf{K}$ полиномиально быстро)}.

     

     \medskip

     

     \noindent

     Д\,о\,к\,а\,з\,а\,т\,е\,л\,ь\,с\,т\,в\,о\,.\ \ Расширенная таблица коммутации $\mathbf{K}^*\hm = 

\{h_1, h_2, \ldots , h_M\}$ порождается из исходной таблицы коммутации добавлением не 

более чем

     $$

     M-n=2s<2n=O(n)

     $$

дополнительных новых строк (размера ровно~$n$ каждая).



     Объектом дальнейшего обсуждения будут <<внутренние взаимосвязи>> как в 

структурах рассматриваемых последовательностей (заголовков пакетов), так и в заранее 

заданных подпоследовательностях символов, требующих проверки на вложимость (в 

такие заголовки). Фактически речь пойдет о множествах \textit{примеров} (конкретных 

экземпляров) заголовков, а также о конечных алфавитах атомарных символов 

(образующих), используемых для <<кодирования>> появления единиц и нулей на 

соответствующих позициях и в заголовках, и в (требующих проверки на вложимость в эти 

заголовки) заданных наборах подпоследовательностей нулей и единиц. Здесь в первую 

очередь будут представлять интерес следующие пары объектов: множество \textit{всех 

примеров}, в которые входят элементы \textit{заданного подмножества образующих}, и 

множество \textit{всех образующих}, одновременно входящих \textit{во все} элементы 

\textit{заданного множества примеров}.

     

     Перейдем к более формальному описанию представленных пар множеств. 

Располагая множествами~$\mathbf{U}^*$ и~$\mathbf{K}^*$, определим два отображения~$f$ и~$\varphi$:

     \begin{gather*}

     \forall \mathbf{A}\in \mbox{\biP}\left(\mathbf{U}^*\right) f(\mathbf{A}) = \{h\in \mathbf{K}^*\ \mbox{таких, что для}\ \forall 

a\in \mathbf{A}\ \mbox{имеет мест}\ a\in \mathbf{h}\}\,;\\

     \forall \mathbf{B}\in \mbox{\biP}\left(\mathbf{K}^*\right) \varphi(\mathbf{B}) = \{a\in \mathbf{U}^*\ \mbox{таких, что для}\ 

\forall h\in \mathbf{B}\ \mbox{имеет мест}\ a\in \mathbf{h}\}\,.

     \end{gather*}

При этом $\mbox{\biP}\left(\mathbf{X}\right)$~--- множество всех подмножеств множества~$\mathbf{X}$ (полагая 

$\mathbf{X}\hm\in \{\mathbf{U}^*,\mathbf{K}^*\}$) . Несложно убедиться, что пара отображений $\langle 

f,\varphi\rangle$ представляет собой соответствие Галуа (см., например,~[8, 9]), а их 

произведения $f(\varphi (\mathbf{B}))$ и $\varphi (f(\mathbf{A}))$~--- соответствующие 

замыкания Галуа~[8, 9]. Через $\mathbf{GC}_{f,\varphi} (\mathbf{K}^*)$ и $\mathbf{GC}_{\varphi,f} 

(\mathbf{U}^*)$ будем обозначать множества неподвижных точек соответствующих замыканий 

Галуа:

\begin{gather*}

\mathbf{GC}_{f,\varphi}(\mathbf{K}^*) =\left\{ \mathbf{B}\in \mbox{\biP}\left(\mathbf{K}^*\right)\ 

\mbox{таких, что}\ 

f(\varphi(\mathbf{B}))=\mathbf{B}\right\}\,;\\

\mathbf{GC}_{\varphi,f}(\mathbf{U}^*) =\left\{ \mathbf{A}\in \mbox{\biP}\left(\mathbf{U}^*\right)\ \mbox{таких, что}\ 

\varphi(f(\mathbf{A}))=\mathbf{A}\right\}\,.

\end{gather*}

Каждое из этих множеств можно рассматривать как частично упорядоченное в 

соответствии как со взаимным вложением соответствующих множеств заголовков (будем 

также называть их объектами), так и подмножеств образующих.



Будем говорить, что слово (а в данном случае~--- и множество букв\footnote{Как уже было показано 

ранее, 

номер позиции каждого символа~$\beta_j$ в этом слове, а также выбор~0 или~1 в качестве 

значения~$\beta_j$ 

может быть закодирован номером соответствующей образующей~--- $a_{0j}$ или~$a_{1j}$. Для кодировки 

выбора третьего значения~${x}$ из множества $\{0,\,1,\,x\}$ в качестве значения для~$\beta_j$ 

используем одновременно пару $a_{0j}, a_{1j}$.})

 $\mathbf{w}_1 \hm= \langle\beta_{11}, \beta_{21}, \ldots , \beta_{p} 

\rangle$ в расширенном алфавите $\mathbf{U}^*\cup \{x\}$ \textit{вкладывается} в 

слово $\mathbf{w}_2 \hm= \langle\beta_{12}, 

\beta_{22}, \ldots , \beta_q\rangle$ (в этом же алфавите) тогда и только тогда, когда одновременно

\begin{enumerate}[(i)]

\item  $ q \geq p$;

\item  для каждого соответствующего\footnote{Напомним, что в каждом из слов~$\mathbf{w}_1$ и $\mathbf{w}_2$ при 

кодировке символами расширенного алфавита $\mathbf{U}^*\cup \{x\}$ используются образующие вида~$a_{rs}$, 

отражающие в том числе позиции соответствующих символов и в слове максимальной длины $2n$.} 

(находящегося на соответствующей позиции) значения $\beta_{j1}$ и~$\beta_{i2}$ из множества $\{0,\,1\}$ в 

каждом из этих слов $\beta_{j1}\hm = \beta_{i2}$ (т.\,е.\ используется один и тот же символ $a_t\in \mathbf{U}^*$);

\item  для соответствующих $\beta_{j1}$ и~$\beta_{i2}$ из $\mathbf{w}_1$ и $\mathbf{w}_2$ если $\beta_{i2}$ 

сопоставляется значению~$x$, то $\beta_{j1}$ сопоставляется одному из 

подходящих $a_t\hm\in U^*$ (для представления значений из множества 

$\{0,\,1\}$) либо подходящей паре вида $a_{0t}, a_{1t}$ (для представления 

значения~$x$).

\end{enumerate}



 Теперь пусть



$\mathbf{h} = \{a_1, a_2, \ldots , a_k\}$~--- заголовок в расширенном алфавите $\mathbf{U}^*\hm = 

\{a_{01}, a_{11}, a_{02}, a_{12}, \ldots , a_{0n}, a_{1n}\}$; 



$\mathbf{K}= \{h_1, h_2, \ldots , h_m\}$~--- таблица коммутации (где $h_i$~--- ее строки, а $i\hm\in\{1, 

2, \ldots , m\}$), переформулированная также в алфавит~$\mathbf{U}^*$; 



$\mathbf{K}^* = \{h_1, h_2, \ldots , h_M\}$~--- таблица коммутации~$\mathbf{K}$, расширенная для каж\-до\-го 

заголовка из множества $\{h_1, h_2, \ldots , h_m\}$, содержащего со\-от\-вет\-ст\-ву\-ющую 

третьему значению~$x$ из множества $\{0,\, 1,\, x\}$ в позиции номер~$i$ пару 

$\{a_{0i}, a_{1i}\}$, дополнительными заголовками, порождаемыми по введенному 

выше {правилу};



$[\mathbf{b}]_{\mathbf{K},\mathbf{U}^*}$~--- определенное ранее замыкание Галуа (сформированное с учетом 

$\mathbf{U}^*$ как исходного алфавита и $\mathbf{K}\hm= \{h_1, h_2, \ldots , h_m\}$ как исходно 

заданного множества объектов) для множества~$\mathbf{b}$; 



$\mathbf{GC}(\mathbf{H},\mathbf{U}^*)$~--- множество всех замыканий Галуа (заданных по 

представленной выше схеме), порожденных на заданном множестве объектов~$\mathbf{H}$, в 

свою очередь сформированных с помощью образующих из множества~$\mathbf{U}^*$.



    \medskip

    

    \noindent

    \textbf{Утверждение 2.} \textit{Заголовок~$\mathbf{h}$ вкладывается в одну из строк 

коммутационной таблицы~$\mathbf{K}$ $($и при этом также содержится в множестве 

$\mathbf{GC}(\mathbf{K}^*, \mathbf{U}^*)$ замыканий Галуа, построенных на расширенной 

коммутационной таблице~$\mathbf{K}^*$$)$ тогда и только тогда, когда}

    \begin{multline}

    \mathop{\bigcup}\limits_{i=1}^k \left[ \{a_{i}\}\right]_{\mathbf{K},\mathbf{U}^*} =

    [\mathbf{h}]_{\mathbf{K},\mathbf{U}^*} 

=[ \{ a_1,a_2,\ldots , a_k\}]_{\mathbf{K},\mathbf{U}^*} =\left[ \mathop{\bigcup}\limits_{i=1}^k 

\{a_i\}\right]_{\mathbf{K},\mathbf{U}^*} ={}\\

{}=\{a_1,a_2, \ldots ,a_k\}\,.

    \label{e1-zab}

    \end{multline}



%\medskip



\noindent

Д\,о\,к\,а\,з\,а\,т\,е\,л\,ь\,с\,т\,в\,о\,.\ \  Прежде всего представим условие доказываемого 

утверждения в более простом виде:

\begin{equation}

\mathbf{h}\in \mathbf{GC}(\mathbf{H},\mathbf{U}^*)\ \mbox{тогда и только, когда}\ 

\mathop{\bigcup}\limits_{i=1}^k \left[\{ a_i\}\right]_{\mathbf{K},\mathbf{U}^*}=\mathbf{h}\,.

\label{e1.1-zab}

\end{equation} 

Далее учтем, что формула $\mathbf{h}\hm\in \mathbf{GC}(\mathbf{H}, \mathbf{U}^*)$ в 

соотношении~(\ref{e1.1-zab}) есть указание на то, что $\mathbf{h}$~--- это неподвижная 

точка замыкания Галуа $[\_ ]_{\mathbf{K},\mathbf{U}^*}$:

$$

[\mathbf{h}]_{\mathbf{K},\mathbf{U}^*} = [\{a_1, a_2, \ldots , a_k\}]_{\mathbf{K},\mathbf{U}^*} = 

\{a_1, a_2, \ldots , a_k\}\,.

$$

Кроме того, по определению замыкания Галуа $[\_ ]_{\mathbf{K},\mathbf{U}^*}$ для всякого $\mathbf{h}_0$ 

имеет место

\begin{equation}

\mathbf{h}\subseteq  [\mathbf{h}]_{\mathbf{K},\mathbf{U}^*}\,.   

\label{e1.2-zab}

\end{equation}

Таким образом, возможны всего два варианта для $[\mathbf{h}]_{\mathbf{K},\mathbf{U}^*}$:

\begin{equation}

\mathbf{h} = \{a_1, a_2, \ldots , a_k\} = \mathop{\bigcup}\limits_{i=1}^k [\{a_i\}]_{\mathbf{K},\mathbf{U}^*}

\label{e1.2.1-zab}

\end{equation}

или

\begin{equation}

\mathbf{h} = \{a_1,a_2, \ldots , a_k\} \subset \mathop{\bigcup}\limits_{i=1}^k 

[\{a_i\}]_{\mathbf{K},\mathbf{U}^*}\,.\label{e1.2.2-zab}

\end{equation}

В случае~(\ref{e1.2.2-zab}) в множестве 

$\mathop{\bigcup}\limits_{i=1}^k [\{a_i\}]_{\mathbf{K},\mathbf{U}^*}$

помимо элементов $a_1, a_2, \ldots , a_k$ содержится также некоторый элемент~$b_0$ из 

$\mathbf{U}^*$, такой что

\begin{equation}

b_0 \not\in  \mathbf{h} = \{a_1, a_2, \ldots , a_k\}\,,

\label{e1.3-zab}

\end{equation}

при этом хотя бы одно из множеств $[\{a_i\}]_{K,U^*}$ ($i \hm= 1, 2, \ldots , k$) также 

содержит этот элемент~$b_0$. Пусть это будет некоторое множество $[\{a_0\}]_{\mathbf{K},\mathbf{U}^*}$. 

Тогда по определению замыкания Галуа $[\_]_{\mathbf{K},\mathbf{U}^*}$ образующая $a_0$ входит в 

соответствующие объекты из $\mathbf{K}$ \textit{только в паре с образующей}~$b_0$. Другими 

словами, так как $[\{a_1, a_2, \ldots , a_k\}]_{\mathbf{K},\mathbf{U}^*}$ есть множество всех образующих из 

$\mathbf{U}^*$, которые одновременно входят в соответствующее подмножество 

\begin{equation}

\mathbf{K}_{a_1, a_2, \ldots  , a_k} \subseteq \mathbf{K}\,,

\label{e1.4-zab}

\end{equation}

то каждый содержащий $a_0$ объект~$\mathbf{h}_0$ из этого подмножества $\mathbf{K}_{a_1, a_2, 

\ldots, a_k}$ содержит также и~$b_0$, т.\,е.\ в этом случае

\begin{equation}

\left(\{a_1, a_2, \ldots , a_k\} \cup \{ b_0\}\right) \subseteq [\{a_1, a_2, \ldots , 

a_k\}]_{\mathbf{K},\mathbf{U}^*}\,,

\label{e1.5-zab}

\end{equation}

откуда следует

\begin{equation}

\{a_1, a_2, \ldots , a_k\} \subset [\{a_1, a_2, \ldots , a_k\}]_{\mathbf{K},\mathbf{U}^*}\,.

\label{e1.6-zab}

\end{equation}

Таким образом, в этом случае $\mathbf{h}\hm = \{a_1, a_2, \ldots , a_k\}$ не является 

неподвижной точкой замыкания Галуа $[\_]_{\mathbf{K},\mathbf{U}^*}$:

\begin{equation}

\mathbf{h}\not\in \mathbf{GC}(\mathbf{H}, \mathbf{U}^*)\,.

\label{e1.7-zab}

\end{equation}

Теперь имеются в наличии все необходимые инструменты, чтобы показать 

справедливость утверждения~(\ref{e1-zab}), переформулированного к 

виду~(\ref{e1.1-zab}).



\medskip



\noindent

\textbf{Необходимость.} Положим, что выполняется формула $\mathbf{h}\hm\in 

\mathbf{GC}(\mathbf{H}, \mathbf{U}^*)$. В~этом случае выполняется лишь одно из 

соотношений~(\ref{e1.2.1-zab}) или (\ref{e1.2.2-zab}) . Пусть это будет~(\ref{e1.2.2-zab}), 

но тогда продвижение по цепочке рассуждений~(\ref{e1.3-zab})--(\ref{e1.7-zab}) приводит 

к противоречию с исходным допущением о том, что $\mathbf{h}\hm\in 

\mathbf{GC}(\mathbf{H},\mathbf{U}^*)$. Таким образом, остается лишь вариант~(\ref{e1.2.1-zab}).

\medskip



\noindent

\textbf{Достаточность.} Положим, что выполнено соотношение

\begin{equation}

\mathop{\bigcup}\limits_{i=1}^k  [\{a_i\}]_{\mathbf{K},\mathbf{U}^*} = \mathbf{h}\,.

\label{e1.8-zab}

\end{equation}

Тогда, предположив, что

\begin{equation}

\mathbf{h}\not\in \mathbf{GC}(\mathbf{H},\mathbf{U}^*)\,,

\label{e1.9-zab}

\end{equation}

и принимая во внимание, что выполнимость~(\ref{e1.9-zab}) сигнализирует о наличии 

среди элементов множества $a_1, a_2, \ldots , a_k$ такого~$a_0$, что его замыкание 

$[\{a_0\}]_{\mathbf{K},\mathbf{U}^*}$ содержит также некоторый элемент~$b_0$ такой, что для него вновь 

выполняется условие~(\ref{e1.3-zab}):

$$

b_0 \not\in \mathbf{h} = \{a_1, a_2, \ldots, a_k\}\,.

$$

Далее по соображениям, учитывавшимся ранее в цепочке 

рассуждений~(\ref{e1.3-zab})--(\ref{e1.7-zab}), несложно прийти к соотношению:

$$

\left(\mathop{\bigcup}\limits_{i=1}^k [\{a_i\}]_{\mathbf{K},\mathbf{U}^*}\right) \supset\mathbf{h}\, ,

$$

которое противоречит исходно сделанному допущению~(\ref{e1.8-zab}).



     В заключение обратим внимание на нетривиальный характер 

соотношения~(\ref{e1.8-zab}), в рамках которого не требуется, чтобы \textit{каждая} их 

входящих в $\mathbf{h}$ образу\-ющих соответствовала бы \textit{неподвижной точке} 

замыкания Галуа $[\_]_{\mathbf{K},\mathbf{U}^*}$. Требуется, чтобы объединение замыканий всех входящих в 

$\mathbf{h}$ образующих $a_1, a_2, \ldots , a_k$ не выводило бы результат за пределы 

множества образующих $\{a_1, a_2, \ldots , a_k\} \hm= \mathbf{h}$.

     

     \medskip

     

     \noindent

     \textbf{Пример.}\ Положим $\mathbf{U}^* \hm= \{a_1, a_2, a_3, a_4, a_5, a_6, a_7\}$, а 

     $\mathbf{K}\hm = \{h_1, h_2, h_3, h_4\}$, где

     \begin{align*}

h_1 &= \{a_1, a_2, a_4\}\,;\\

h_2 &= \{a_1, a_2, a_5\}\,;\\

h_3 &= \{a_1, a_3, a_6\}\,;\\

h_4 &= \{a_1, a_3, a_7\}\,.

\end{align*}

Несложно убедиться, что

\begin{align*}

[\{a_1\}]_{\mathbf{K},\mathbf{U}^*}& = \{a_1\}\,;\\

[\{a_2\}]_{\mathbf{K},\mathbf{U}^*} & = \{a_1, a_2\} \supset\{a_2\}\,;\\

[\{a_3\}]_{\mathbf{K},\mathbf{U}^*}& = \{a_1, a_3\} \supset\{a_3\}\,;

\end{align*}

тем не менее,

\begin{align*}

h_1 &= [\{a_1, a_2, a_4\}]_{\mathbf{K},\mathbf{U}^*} = [\{a_4\}]_{\mathbf{K},\mathbf{U}^*} = \{a_1, a_2, a_4\}\supset \{a_4\}\,;\\

h_2 &= [\{a_1, a_2, a_5\}]_{\mathbf{K},\mathbf{U}^*} = [\{a_5\}]_{\mathbf{K},\mathbf{U}^*} = \{a_1, a_2, a_5\} \supset\{a_5\}\,;\\

h_3 &= [\{a_1, a_3, a_6\}]_{\mathbf{K},\mathbf{U}^*} = [\{a_6\}]_{\mathbf{K},\mathbf{U}^*} = \{a_1, a_3, a_5\} \supset\{a_6\}\,;\\

h_4 &= [\{a_1, a_3, a_7\}]_{\mathbf{K},\mathbf{U}^*} = [\{a_7\}]_{\mathbf{K},\mathbf{U}^*} = \{a_1, a_3, a_7\} \supset\{a_7\}\,.

\end{align*}



Таким образом, для проверки присутствия заголовка~$\mathbf{h}$ в коммутационной 

таблице~$\mathbf{K}$ следует убедиться, что~$\mathbf{h}$ есть неподвижная точка 

соответствующего замыкания Галуа $[\_]_{\mathbf{K},\mathbf{U}^*}$, т.\,е.\ 

$[\mathbf{h}]_{\mathbf{K},\mathbf{U}^*}\hm = [\{a_1, 

a_2, \ldots , a_k\}]_{\mathbf{K},\mathbf{U}^*}\hm = [\mathbf{h}]_{\mathbf{K},\mathbf{U}^*}

\hm= \{a_1, a_2, \ldots , a_k\}$. При этом 

алгоритмически это можно сделать, проверив, например, выполнение условия:

\begin{equation}

\{a_1, a_2, \ldots , a_k\} =\mathop{\bigcup}\limits_{i=1}^k [\{a_i\}]_{\mathbf{K},\mathbf{U}^*}\,.

\label{e2-zab}

\end{equation}

     

     \medskip

     

     \noindent

     \textbf{Замечание 1.} Соотношение~(\ref{e2-zab}) при диагностике встречаемости 

попадающих на вход заголовков~$\mathbf{h}$ в коммутационной таблице~$\mathbf{K}$ позволяет:

     \begin{itemize}

\item один раз рассчитать необходимые значения для замыканий (в смысле текущего 

состояния коммутационной таблицы~$\mathbf{K}$) всех элементов алфавита~$\mathbf{U}^*$;

\item использовать эти расчеты многократно (в случае отсутствия изменений в 

таб\-ли\-це~$\mathbf{K}$) для диагностики встречаемости в ней заголовков вновь приходящих 

на вход пакетов.

\end{itemize}



     \noindent

     \textbf{Замечание 2.} Соотношение~(\ref{e2-zab}) прямо не зависит от 

параметра~$m$~--- числа строк в таблице коммутации~$\mathbf{K}$, что позволяет проводить 

(после однократного вычисления замыканий для элементов множества~$\mathbf{U}^*$) 

диагностику встречаемости вновь попадающих на вход заголовков~$\mathbf{h}$ за не 

зависящее от размеров таб\-ли\-цы~$\mathbf{K}$ время $T\hm = const (m)$.

     

     При этом размер (число строк) таблицы~$\mathbf{K}$ существен лишь при вычислении 

замыканий образующих из множества~$\mathbf{U}^*$ с учетом объектов из множества $\mathbf{K}\hm = 

\{h_1, h_2, \ldots , h_m\}$, в которые эти образующие входят.

     



\section{Задача о вложимости слов и~алгоритм ее~решения}



     Теперь обратимся собственно к задаче о проверке вложимости заданного набора слов 

в текущее слово (например, в заголовок пакета, поступающего на вход сетевого 

устройства). Пусть $\mathbf{P}\hm = \{\mbox{\bip}_1, \mbox{\bip}_2, \ldots , \mbox{\bip}_r\}$~--- некоторое множество слов вида 

$\mbox{\bip}_i \hm= \langle \alpha_{1i}, \alpha_{2i}, \ldots , \alpha_{k(i)}\rangle$, каждое соответствующей 

длины $k(i) \hm< n$, построенных лишь с помощью символов~0 и~1.

     

    Нормализованным видом каждого $\mbox{\bip}_i$ назовем его расширение $\mathbf{p}_i$ в 

алфавите~$\mathbf{U}^*$ до слова длины $l \hm= k(i) \hm+ 2(n \hm- k(i)) \hm= 2n \hm- k(i)$, в 

котором оставшиеся до <<стандартного>> размера заголовков~$\mathbf{h}$ 

(<<незанятые>>) позиции сперва дополнены значениями~$x$, каждое из которых далее (в 

нормализованной записи для каждго слова~$\mbox{\bip}_i$) заменено на соответствующую пару 

символов вида $\{a_{0i}, a_{1i}\}$ из~$\mathbf{U}^*$. Таким образом множество слов~$\mathbf{P}$ 

преобразуется в нормализованное множество слов $\mathbf{P}\hm = \{\mathbf{p}_1, 

\mathbf{p}_2, \ldots , \mathbf{p}_r\}$ в алфавите~$\mathbf{U}^*$.

    

    Построим по $\mathbf{P}\hm = \{\mathbf{p}_1, \mathbf{p}_2, \ldots , \mathbf{p}_r\}$ с 

применением установленного выше {правила} расширенное множество 

$\mathbf{P}^*\hm = \{\mathbf{p}_1, \mathbf{p}_2, \ldots , \mathbf{p}_R\}$. 

    

    \medskip

    

    \noindent

    \textbf{Утверждение 3.} \textit{Множество $\mathbf{P}^*$ по множеству~$\mathbf{P}$ 

строится полиномиально быстро. Это очевидное следствие утверждения}~1.

    

    \medskip

    

    \noindent

    \textbf{Утверждение 4.} \textit{Для каждого исходно заданного слова $\mbox{\bip}_i \hm= 

\langle\alpha_{1i}, \alpha_{2i}, \ldots , \alpha_{k(i)}\rangle$ из множества $\mathbf{P}\hm = \{\mbox{\bip}_1, 

\mbox{\bip}_2, \ldots 

, \mbox{\bip}_r\}$ в множестве всех замыканий Галуа $\mathbf{GC}(\mathbf{P}^*,\mathbf{U}^*)$, 

построенных на объектах из $\mathbf{P}^* \hm= \{\mathbf{p}_1, \mathbf{p}_2, \ldots , 

\mathbf{p}_R\}$ найдутся все его возможные <<дополнения>> $($в смысле расстановки~$0$ 

и~$1$ в <<незанятые>> символами $\alpha_{1i}, \alpha_{2i}, \ldots , \alpha_{k(i)}$ позиции 

<<объемлющего>> слово $\mbox{\bip}_i$ булевского вектора длины~$n)$}.



\medskip



\noindent

Д\,о\,к\,а\,з\,а\,т\,е\,л\,ь\,с\,т\,в\,о\,.\ \ Выберем из множества $\mathbf{P}\hm = 

\{\mbox{\bip}_1, \mbox{\bip}_2, \ldots , 

\mbox{\bip}_r\}$ произвольный элемент~$\mbox{\bip}_0$. Для каждого дополняющего $\mbox{\bip}_0$ до 

<<канонической>> длины $n$ значения~${x}$ в позиции~$i_0$ воспользуемся объектами 

$\mathbf{p}_0$ и $\mathbf{p}_0^{i0}$, а также объектами $\mathbf{p}_0$ и $\mathbf{p}_0^{i1}$ 

(см.\ таблицу), чтобы получить с их помощью оба <<продолжения>> слова $\mbox{\bip}_0$, 

имеющих соответственно~0 и~1 в позиции~$i_0$. Проделав эту операцию для каждой из 

соответствующих $x$-позиций (<<вокруг>> исходных образующих слова~$\mbox{\bip}_0$) и 

воспользовавшись полученными <<покоординатными>> заменами значений~$x$ на 0 и~1 

для удаления всех имевшихся $x$-зна\-че\-ний из соответствующих расширений 

слова~$\mbox{\bip}_0$ до булевского вектора длины~$n$, получаем (уже в рамках множества 

$\mathbf{GC}(\mathbf{P}^*,\mathbf{U}^*)$) каждое из искомых расширений.





Теперь для проверки вложимости слов из заданного множества $\mathbf{P}\hm= \{

\mbox{\bip}_1, \mbox{\bip}_2, \ldots , 

\mbox{\bip}_r\}$ в заданный булевский вектор $\mathbf{h} \hm= \langle \alpha_1, \alpha_2, \ldots , \alpha_n\rangle$ 

длины~$n$ достаточно (следуя Утверждению~2) реализовать следующий алгоритм:

\begin{enumerate}[(1)]

\item построить на объектах множества $\mathbf{P}^*$ замыкания Галуа для каждой из 

образующих из множества~$\mathbf{U}^*$;

\item перекодировать $\mathbf{h}$ в алфавит~$\mathbf{U}^*$ по уже обсуждавшимся правилам 

(превратив его в соответствующее множество образующих из~$\mathbf{U}^*$);

\item проверить выполнение условия~(\ref{e2-zab}) как следствия утверждения~2;

\item выделить (идентифицировав вложимость множества~$\mathbf{h}$, 

порожденного по~$\mbox{\bih}$ перекодировкой в алфавит~$\mathbf{U}^*$, в соответствующие 

строки таблицы~$\mathbf{P}^*$) вкладывающиеся в~$\mbox{\bih}$ слова из множества~$\mathbf{P}$.

\end{enumerate}



При этом (на основании Замечаний~1 и~2) процедура поиска ответа на по\-став\-лен\-ный 

вопрос о вложимости подслов из множества~$\mathbf{P}$ в текущее слово (заголовок)~$\mbox{\bih}$ после 

однократного порождения соответствующих замыканий Галуа на всех образующих из 

множества~$\mathbf{U}^*$ может исполняться за время, не зависящее от размеров множества~$\mathbf{P}$.





\section{Заключение}



     Таким образом, при использовании ПКС-технологий имеется возможность, вынося 

вычисления соответствующих замыканий Галуа на внешний ПКС-конт\-рол\-лер, а также 

используя ТСАМ (Ternary Content Addressible Memory)

для проверки условия~(\ref{e2-zab}) как следствия утверждения~2 

(позволяющего убедиться, что вы действительно имеете дело с соответствующей 

неподвижной точкой построенного замыкания Галуа), определенным образом сократить 

время обработки входных сообщений на используемых сетевых устройствах. 

Определенные основания ожидать ускорения работы сети в рамках предлагаемого 

подхода дает \textit{однократное} (при фиксированной структуре используемой таблицы 

коммутации) построение \textbf{GC}-<<ба\-зи\-са>>~--- множества замыканий 

одноэлементных подмножеств образующих из алфавита $\mathbf{U}^*$, \textit{многократно} 

используемого затем для <<диагностики>> заголовков вновь приходящих входных 

пакетов сообщений. Наконец, имеющиеся здесь также возможности использовать на 

соответствующем ПКС-конт\-рол\-ле\-ре \textit{параллельные вычисления} (в частности, 

при анализе отнесения входных пакетов к разным портам коммутатора или же при 

вычислении замыканий Галуа для каждого из элементов множества образующих~$\mathbf{U}^*$) 

позволяют надеяться на определенное дополнительное ускорение процесса обработки 

сообщений в компьютерных сетях на базе ПКС-тех\-но\-логий.



 {\small\frenchspacing

{%\baselineskip=10.8pt

\addcontentsline{toc}{section}{Литература}

\begin{thebibliography}{9}



\bibitem{1-zab}

\Au{Смелянский Р.\,Л.} Компьютерные сети: В 2-х т.~--- М: Академия, 2011.

\bibitem{2-zab}

\Au{Гэри М., Джонсон Д.\,С.} Вычислительные машины и труднорешаемые задачи.~--- 

М.: Мир, 1982. 416~с.

\bibitem{3-zab}

Open Networking Foundation. {\sf https://www.opennetworking.org}; {\sf 

http://www.\linebreak openflow.org}. 

\bibitem{4-zab}

Open Networking Laboratory$\backslash$Open Networking Research Center. {\sf 

http://onrc.\linebreak stanford.edu}.

\bibitem{5-zab}

Global Environment for Network Innovations (GENI): Национальная программа США. {\sf 

http://www.geni.net}.

\bibitem{6-zab}

OFELIA: Исследовательский проект 7-й Рамочной программы Европейского Союза. {\sf 

http://www.fp7-ofelia.eu}.

\bibitem{7-zab}

FEDERICA: Исследовательский проект 7-й Рамочной программы Европейского Союза. 

{\sf http://dana.i2cat.net/the-end-of-federica-project/platform}.



\bibitem{9-zab}

\Au{Кон П.} Универсальная алгебра.~--- М.: Мир, 1968. 359~с.



 \label{end\stat}

\bibitem{8-zab}

\Au{Гусакова С.\,М., Финн В.\,К.} Сходства и правдоподобный вывод~/\!\!/ Известия АН 

СССР. Сер. Техническая кибернетика, 1987. №\,5. С.~42--63. 





 



 \end{thebibliography}

}

}

